\section{Discussion on Adaptive Speculative Strategies Derived from Empirical Findings}
\label{sec:acsd}
The current empirical study identifies the strongest fixed-policy speedup at
a smaller block size ($k{=}4$), with net speed degrading as $k$ increases, even
when the accepted block length rises. This pattern indicates that fixed-policy
drafting is vulnerable to acceptance/overhead imbalance and long-tail latency.
We therefore design and evaluate an adaptive-policy variant that keeps standard
target-side verification while adapting control decisions online.

ACSD is intentionally not framed as a new peak-throughput optimizer. Its role
is to provide a practical tail-latency guardrail under dynamic acceptance
collapse, while keeping the mean performance near the best static operating point.

Advanced variants such as EAGLE-3 and DFlash reduce drafting cost through new
drafter architectures and dedicated training pipelines, often evaluated on
strong accelerators. Our objective is different: to maintain a stack of drafters / verifiers of the same family
 that is deployable on consumer hardware, and recover
speed and stability by runtime control (adaptive $k$, rescue switching, and
tail fallback) rather than a newly trained drafter.

\subsection{Design objective}

Let speculative-cycle latency be
\begin{equation}
  L = \frac{T_{\text{draft}} + T_{\text{verify}}}{\tau},
\end{equation}
where $\tau$ is the expected number of accepted tokens per cycle. In fixed-policy speculative
decoding, poorly matched samples can reduce $\tau$ and increase total verify
steps, producing a heavy-tail runtime. ACSD explicitly targets this effect by
adapting draft policy at the sample level.


\subsection{Formalization of the ACSD Policy}
To rigorously define the runtime behavior, we formalize the ACSD controller as a deterministic finite-state machine (FSM). Let $t$ denote the current verification step, and $k_t$ be the dynamic draft length. The system state is defined as $S_t \in \{\mathcal{D}_{pri}, \mathcal{D}_{res}, \mathcal{M}_{AR}\}$, representing the primary drafter (0.5B), the rescue drafter (1.5B), and the autoregressive fallback mode, respectively. The state transitions and scheduling parameters are directly governed by the real-time empirical acceptance rate $\alpha_t = \frac{\tau_t}{k_t}$, where $\tau_t$ is the number of accepted tokens at step $t$. ACSD incorporates four core online control mechanisms:

\begin{enumerate}
    \item \textbf{Dynamic Draft Length ($k_t$):} While operating in $\mathcal{D}_{pri}$ or $\mathcal{D}_{res}$, $k_t$ is modulated within a narrow operational neighborhood around the empirical optimum of 4 (typically 3 to 5). Specifically, the controller updates the length via $k_{t+1} = \min(k_t + 1, k_{max})$ if $\alpha_t \geq \theta_{up}$, and $k_{t+1} = \max(k_t - 1, k_{min})$ if $\alpha_t < \theta_{down}$. This bounding strategy effectively caps the memory bandwidth overhead during target-side verification.

    \item \textbf{Cascaded Draft Switch ($S_t \rightarrow \mathcal{D}_{res}$):} We maintain a continuous penalty counter $C_t$ to track consecutive low-acceptance events. If $\alpha_t < \theta_{rescue}$, the counter increments; otherwise, $C_t = 0$. Once $C_t \geq N_{rescue}$, the controller triggers a state transition to the 1.5B rescue drafter ($S_{t+1} = \mathcal{D}_{res}$) to leverage its higher representational capacity for difficult distributions. To mitigate scheduling oscillation, a cooling window $W_{cool}$ (measured in generation steps) is strictly enforced before any subsequent model switching.

    \item \textbf{Tail Latency Guardrail ($S_t \rightarrow \mathcal{M}_{AR}$):} To prevent severe tail-latency degradation under sustained high-entropy inputs, a terminal fallback mechanism is employed. If the rescue drafter fails to recover the acceptance rate (i.e., $\alpha_t < \theta_{fallback}$ for $N_{fallback}$ consecutive steps) after generating a minimum prefix length $L_{min}$, the system forcefully transitions to pure autoregressive decoding ($S_{t+1} = \mathcal{M}_{AR}$) for the remainder of the sample.

    \item \textbf{Checkpointed Execution:} The execution context, including partial generation sequences and verification logs, is serialized at predefined intervals to guarantee system resilience and allow safe resumption of interrupted runs.
\end{enumerate}

\subsection{How this guidance is used}

Future adaptive-policy evaluation should remain a strict incremental study under
the same protocol, with primary comparison against fixed $k{=}4$ and high-$k$
($k{=}16$) operating points.

\subsection{Observed ACSD outcomes}

Measured ACSD outcomes are reported in Section~\ref{sec:results}
(Table~\ref{tab:acsd-headline}). In brief, ACSD achieves strong acceleration
relative to autoregressive baselines in both regimes
($S{=}2.9561$ deterministic, $S{=}2.8658$ stochastic; both $p<10^{-4}$), while
keeping fallback rates low (1.6\% and 2.2\%). However, ACSD does not exceed
the fastest fixed-policy setting (0.5B, $k{=}4$ deterministic at
$S{=}3.0674$).

For deployment, this should be read as a stability trade: ACSD sacrifices a
small amount of mean throughput versus fixed $k{=}4$, but keeps tails close to
that operating point and far from the high-$k$ tail-latency regime
(Table~\ref{tab:acsd-tail}), while using rescue/fallback only when needed.

\subsection{Scope in this paper}

This section combines design rationale with measured ACSD behavior under the
same protocol. Quantitative claims are restricted to the reported fixed policy
and ACSD runs in this manuscript; broader generalization still requires
additional hardware/model families.
